Het uitspansel vibreerde

Het engelennummer **1711** is een krachtige combinatie van de energieën en vibraties van de nummers 1 en 7. Hier is een interpretatie van de betekenis:

1. **Nummer 1** komt twee keer voor, wat de energie versterkt. Het staat voor:
- Nieuwe beginnen
- Creativiteit
- Doorzettingsvermogen
- Het creëren van je eigen realiteit met positieve gedachten en intenties

2. **Nummer 7** resoneert met:
- Spirituele ontwaking en ontwikkeling
- Innerlijke wijsheid
- Intuïtie en introspectie
- Een dieper begrip van jezelf en je levensdoel

3. **Samengevoegd in 1711**:
- Het getal moedigt je aan om te vertrouwen op je innerlijke wijsheid en intuïtie terwijl je nieuwe stappen in je leven zet.
- Het herinnert je eraan dat je spiritueel wordt geleid, en dat je op de goede weg bent naar het vervullen van je levensmissie.
- Het is ook een aansporing om vol vertrouwen te blijven, ondanks eventuele uitdagingen, omdat deze bijdragen aan je groei en spirituele ontwikkeling.

Als je **1711** vaak tegenkomt, is dat een boodschap van de engelen om je te verzekeren dat je op een positieve, vruchtbare koers zit. Je inspanningen en focus worden ondersteund door het universum, en je wordt aangemoedigd om je unieke gaven en talenten in te zetten voor zowel je eigen welzijn als dat van anderen.

Emmy had opgegeven
met een opgeheven hoofd.

Turend in de duisternis
naar Martino’s schaduwen,
vast in de gevangenis
van Emmy’s oogzenuwen,

zinkend in een diepe trance,
afdalend in een afgrond,
deed Emmy een werveldans
tot ze op haar benen stond.

Daar viel over haar gemoed
een theatrale nevel,
die haar minder fel begroet
in emoties als knevel.

Martino’s brute daden,
werden gedempte echo's,
die al flikkerend baadden,
kleurloos, geurloos en smaakloos

en enkel van korte duur,
op de wanden van een grot.
De figuren op de muur
speelden wel hun duister plot,

maar waren louter schimmen.
Emmy steeg er bovenuit
en liet het licht hen dimmen,
dat straalde van bovenuit.

Deze holle projecties
waren in hun wezen leeg.
Zijn lusten en agressies
een koekje van eigen deeg,

waarop je kon uitzoomen
en niet hoefde te volgen.
Ze konden niet verbloemen,
dat ze niet voor je zorgen,

maar een leugen vertelden.
Over lust als oplossing.
Wat ze op de mouw spelden,
was hun weg naar verlichting.

Ze zette ze op afstand
en ontving de helderheid
dat dit lijf werd aangerand,
maar niet haar feitelijkheid.

Want zij is de toeschouwer
in haar levenstheater.
Wie was de decorbouwer
en wie was de uitbater?

Zoals Topol Tevye
vertolkte als violist
op het dak van de vrije
woning van een scenarist.

Je opmerking brengt een klassiek dilemma naar voren: angst en passie creëren een spanningsveld tussen veiligheid en vrijheid. Dit conflict is als thema door de menselijke geschiedenis heen zichtbaar geweest, zowel in filosofie als in literatuur en psychologie.

Analyse vanuit literair oogpunt:
- **Logica:** Angst staat vaak aan de kant van veiligheid—je kiest wat vertrouwd is, beschermt jezelf tegen het onbekende. Passie daarentegen duwt richting vrijheid: het verlangen, het avontuur, het onbekende.
- **Boek- en historische correctheid:** In talloze werken (denk aan Goethe, Freud, Camus) wordt deze dynamiek uitgewerkt: het is een universeel menselijke vraag—meegaan met je passie betekent risico’s nemen, terwijl veiligheid vaak het onderdrukken van diepere verlangens inhoudt.
- **Rijmschema en lettergrepen:** Er zijn geen rijmregels in je zin, en ook geen vastgesteld metrum. Wil je dit thema in poëtische vorm verwerken (met bijvoorbeeld een ABAB-schema), dan kan ik je hierbij helpen.

Indien gewenst, kan ik dit thema uitwerken in een plot, een canto, of een liedtekst. Geef gerust aan hoe je het conflict tussen angst, passie, veiligheid en vrijheid verder wilt verkennen binnen de context van Emmy Nativus.
